\documentclass[UTF8,12pt]{ctexart}
\usepackage[a4paper,margin=24mm]{geometry}
\usepackage{amsmath,amssymb,bm,graphicx,adjustbox}
\newcommand{\fitmath}[1]{\begin{adjustbox}{max width=\linewidth}$\displaystyle #1$\end{adjustbox}}
\setlength{\parindent}{0pt}
\setlength{\parskip}{.7em}
\sloppy
\allowdisplaybreaks
\begin{document}
\section*{2006 年考研数学三 · 真题}
\subsection*{2006年全国硕士研究生招生考试试题}

\subsection*{一、填空题（本题共6小题，每小题4分，满分24分）}

（1）\(\displaystyle \lim_{n\to\infty}\left(\dfrac{\displaystyle n+1}{\displaystyle n}\right)^{(-1)^n}=\)\_\_\_\_\_\_。


（2）设函数 \(\displaystyle f(x)\) 在 \(\displaystyle x=2\) 的某邻域内可导，且 \(\displaystyle f^{\prime}(x)=e^{f(x)}\)，\(\displaystyle f(2)=1\)，则 \(\displaystyle f^{\prime\prime\prime}(2)=\)\_\_\_\_\_\_。


（3）设函数 \(\displaystyle f(u)\) 可微，且 \(\displaystyle f^{\prime}(0)=\dfrac{\displaystyle 1}{\displaystyle 2}\)，则 \(\displaystyle z=f(4x^2-y^2)\) 在点 \(\displaystyle (1,\allowbreak 2)\) 处的全微分 \(\displaystyle \left.\mathrm dz\right|_{(1,\allowbreak 2)}=\)\_\_\_\_\_\_。


（4）设矩阵 \(\displaystyle A=\begin{pmatrix}2&1\\-1&2\end{pmatrix}\)，\(\displaystyle E\) 为2阶单位矩阵，矩阵 \(\displaystyle B\) 满足 \(\displaystyle BA=B+2E\)，则 \(\displaystyle |B|=\)\_\_\_\_\_\_。


（5）设随机变量 \(\displaystyle X\) 与 \(\displaystyle Y\) 相互独立，且均服从区间 \(\displaystyle [0,\allowbreak 3]\) 上的均匀分布，则 \(\displaystyle P\{\displaystyle \max\{X,\allowbreak Y\}\le1\}=\)\_\_\_\_\_\_。


（6）设总体 \(\displaystyle X\) 的概率密度为 \(\displaystyle f(x)=\dfrac{\displaystyle 1}{\displaystyle 2}e^{-|x|}\;(-\infty<x<+\infty)\)，\(\displaystyle X_1,\allowbreak X_2,\allowbreak \ldots,\allowbreak X_n\) 为总体 \(\displaystyle X\) 的简单随机样本，其样本方差为 \(\displaystyle S^2\)，则 \(\displaystyle E(S^2)=\)\_\_\_\_\_\_。


\subsection*{二、选择题（本题共8小题，每小题4分，满分32分）}

（7）设函数 \(\displaystyle y=f(x)\) 具有二阶导数，且 \(\displaystyle f^{\prime}(x)>0\)，\(\displaystyle f^{\prime\prime}(x)>0\)。\(\displaystyle \Delta x\) 为自变量 \(\displaystyle x\) 在点 \(\displaystyle x_0\) 处的增量，\(\displaystyle \Delta y\) 与 \(\displaystyle \mathrm dy\) 分别为 \(\displaystyle f(x)\) 在点 \(\displaystyle x_0\) 处对应的增量与微分，若 \(\displaystyle \Delta x>0\)，则（　）。（A）\(\displaystyle 0<\mathrm dy<\Delta y\)。（B）\(\displaystyle 0<\Delta y<\mathrm dy\)。（C）\(\displaystyle \Delta y<\mathrm dy<0\)。（D）\(\displaystyle \mathrm dy<\Delta y<0\)。


（8）设函数 \(\displaystyle f(x)\) 在 \(\displaystyle x=0\) 处连续，且 \(\displaystyle \lim_{h\to0}\dfrac{\displaystyle f(h^2)}{\displaystyle h^2}=1\)，则（　）。（A）\(\displaystyle f(0)=0\) 且 \(\displaystyle f^{\prime}_{-}(0)\) 存在。（B）\(\displaystyle f(0)=1\) 且 \(\displaystyle f^{\prime}_{-}(0)\) 存在。（C）\(\displaystyle f(0)=0\) 且 \(\displaystyle f^{\prime}_{+}(0)\) 存在。（D）\(\displaystyle f(0)=1\) 且 \(\displaystyle f^{\prime}_{+}(0)\) 存在。


（9）若级数 \(\displaystyle \sum_{n=1}^{\infty}a_n\) 收敛，则级数（　）。（A）\(\displaystyle \sum|a_n|\) 收敛。（B）\(\displaystyle \sum(-1)^na_n\) 收敛。（C）\(\displaystyle \sum a_na_{n+1}\) 收敛。（D）\(\displaystyle \sum\dfrac{\displaystyle a_n+a_{n+1}}{\displaystyle 2}\) 收敛。


（10）设非齐次线性微分方程 \(\displaystyle y^{\prime}+P(x)y=Q(x)\) 有两个不同的解 \(\displaystyle y_1(x),\allowbreak y_2(x)\)，\(\displaystyle C\) 为任意常数，则该方程的通解是（　）。（A）\(\displaystyle C[y_1(x)-y_2(x)]\)。（B）\(\displaystyle y_1(x)+C[y_1(x)-y_2(x)]\)。（C）\(\displaystyle C[y_1(x)+y_2(x)]\)。（D）\(\displaystyle y_1(x)+C[y_1(x)+y_2(x)]\)。


（11）设 \(\displaystyle f(x,\allowbreak y)\) 与 \(\displaystyle \varphi(x,\allowbreak y)\) 均为可微函数，且 \(\displaystyle \varphi_y^{\prime}(x,\allowbreak y)\ne0\)，已知 \(\displaystyle (x_0,\allowbreak y_0)\) 是 \(\displaystyle f(x,\allowbreak y)\) 在约束条件 \(\displaystyle \varphi(x,\allowbreak y)=0\) 下的一个极值点，下列选项正确的是（　）。（A）若 \(\displaystyle f_x^{\prime}(x_0,\allowbreak y_0)=0\)，则 \(\displaystyle f_y^{\prime}(x_0,\allowbreak y_0)=0\)。（B）若 \(\displaystyle f_x^{\prime}(x_0,\allowbreak y_0)=0\)，则 \(\displaystyle f_y^{\prime}(x_0,\allowbreak y_0)\ne0\)。（C）若 \(\displaystyle f_x^{\prime}(x_0,\allowbreak y_0)\ne0\)，则 \(\displaystyle f_y^{\prime}(x_0,\allowbreak y_0)=0\)。（D）若 \(\displaystyle f_x^{\prime}(x_0,\allowbreak y_0)\ne0\)，则 \(\displaystyle f_y^{\prime}(x_0,\allowbreak y_0)\ne0\)。


（12）设 \(\displaystyle \alpha_1,\allowbreak \ldots,\allowbreak \alpha_s\) 均为 \(\displaystyle n\) 维列向量，\(\displaystyle A\) 是 \(\displaystyle m\times n\) 矩阵，下列选项正确的是（　）。（A）若 \(\displaystyle \alpha_i\) 线性相关，则 \(\displaystyle A\alpha_i\) 线性相关。（B）若 \(\displaystyle \alpha_i\) 线性相关，则 \(\displaystyle A\alpha_i\) 线性无关。（C）若 \(\displaystyle \alpha_i\) 线性无关，则 \(\displaystyle A\alpha_i\) 线性相关。（D）若 \(\displaystyle \alpha_i\) 线性无关，则 \(\displaystyle A\alpha_i\) 线性无关。


（13）设 \(\displaystyle A\) 为3阶矩阵，将 \(\displaystyle A\) 的第2行加到第1行得 \(\displaystyle B\)，再将 \(\displaystyle B\) 的第1列的 \(\displaystyle -1\) 倍加到第2列得 \(\displaystyle C\)，记 \(\displaystyle P=\begin{pmatrix}1&1&0\\0&1&0\\0&0&1\end{pmatrix}\)，则（　）。（A）\(\displaystyle C=P^{-1}AP\)。（B）\(\displaystyle C=PAP^{-1}\)。（C）\(\displaystyle C=P^{\mathrm T}AP\)。（D）\(\displaystyle C=PAP^{\mathrm T}\)。


（14）设随机变量 \(\displaystyle X\) 服从正态分布 \(\displaystyle N(\mu_1,\allowbreak \sigma_1^2)\)，随机变量 \(\displaystyle Y\) 服从正态分布 \(\displaystyle N(\mu_2,\allowbreak \sigma_2^2)\)，且 \(\displaystyle P\{|X-\mu_1|<1\}>P\{|Y-\mu_2|<1\}\)，则必有（　）。（A）\(\displaystyle \sigma_1<\sigma_2\)。（B）\(\displaystyle \sigma_1>\sigma_2\)。（C）\(\displaystyle \mu_1<\mu_2\)。（D）\(\displaystyle \mu_1>\mu_2\)。


\subsection*{三、解答题（本题共9小题，满分94分。解答应写出文字说明、证明过程或演算步骤）}

（15）（本题满分7分）设 \(\displaystyle f(x,\allowbreak y)=\dfrac{\displaystyle y}{\displaystyle 1+xy}-\dfrac{\displaystyle 1-y\sin\dfrac{\displaystyle \pi x}{\displaystyle y}}{\displaystyle \arctan x}\)，\(\displaystyle x>0,\allowbreak y>0\)。求：（I）\(\displaystyle g(x)=\displaystyle \lim_{y\to+\infty}f(x,\allowbreak y)\)；（II）\(\displaystyle \lim_{x\to0^+}g(x)\)。


（16）（本题满分7分）计算二重积分 \(\displaystyle \iint_D\sqrt{y^2-xy}\,\mathrm dx\mathrm dy\)，其中 \(\displaystyle D\) 是由直线 \(\displaystyle y=x,\allowbreak y=1,\allowbreak x=0\) 所围成的平面区域。


（17）（本题满分10分）证明：当 \(\displaystyle 0<a<b<\pi\) 时，\(\displaystyle b\sin b+2\cos b+\pi b>a\sin a+2\cos a+\pi a\)。


（18）（本题满分8分）在 \(\displaystyle xOy\) 坐标平面上，连续曲线 \(\displaystyle L\) 过点 \(\displaystyle M(1,\allowbreak 0)\)，其上任意点 \(\displaystyle P(x,\allowbreak y)\;(x\ne0)\) 处的切线斜率与直线 \(\displaystyle OP\) 的斜率之差等于 \(\displaystyle ax\)（常数 \(\displaystyle a>0\)）。（I）求 \(\displaystyle L\) 的方程；（II）当 \(\displaystyle L\) 与直线 \(\displaystyle y=ax\) 所围成平面图形的面积为 \(\displaystyle \dfrac{\displaystyle 8}{\displaystyle 3}\) 时，确定 \(\displaystyle a\) 的值。


（19）（本题满分10分）求幂级数 \(\displaystyle \sum_{n=1}^{\infty}\dfrac{\displaystyle (-1)^{n-1}x^{2n+1}}{\displaystyle n(2n-1)}\) 的收敛域及和函数 \(\displaystyle S(x)\)。


（20）（本题满分13分）设4维向量组 \(\displaystyle \alpha_1=(1+a,\allowbreak 1,\allowbreak 1,\allowbreak 1)^{\mathrm T}\)，\(\displaystyle \alpha_2=(2,\allowbreak 2+a,\allowbreak 2,\allowbreak 2)^{\mathrm T}\)，\(\displaystyle \alpha_3=(3,\allowbreak 3,\allowbreak 3+a,\allowbreak 3)^{\mathrm T}\)，\(\displaystyle \alpha_4=(4,\allowbreak 4,\allowbreak 4,\allowbreak 4+a)^{\mathrm T}\)。问 \(\displaystyle a\) 为何值时，\(\displaystyle \alpha_1,\allowbreak \alpha_2,\allowbreak \alpha_3,\allowbreak \alpha_4\) 线性相关？当线性相关时，求出一个极大线性无关组，并将其余向量用该极大线性无关组线性表示。


（21）（本题满分13分）设3阶实对称矩阵 \(\displaystyle A\) 的各行元素之和均为3，向量 \(\displaystyle \alpha_1=(-1,\allowbreak 2,\allowbreak -1)^{\mathrm T}\)，\(\displaystyle \alpha_2=(0,\allowbreak -1,\allowbreak 1)^{\mathrm T}\) 是线性方程组 \(\displaystyle Ax=0\) 的两个解。


（21）（续）（I）求 \(\displaystyle A\) 的特征值与特征向量；（II）求正交矩阵 \(\displaystyle Q\) 和对角矩阵 \(\displaystyle \Lambda\)，使得 \(\displaystyle Q^{\mathrm T}AQ=\Lambda\)；（III）求 \(\displaystyle A\) 及 \(\displaystyle (A-\dfrac{\displaystyle 3}{\displaystyle 2}E)^6\)，其中 \(\displaystyle E\) 为3阶单位矩阵。


（22）（本题满分13分）设随机变量 \(\displaystyle X\) 的概率密度为 \(\displaystyle f_X(x)=\begin{cases}\dfrac{\displaystyle 1}{\displaystyle 2},\allowbreak &-1<x<0,\allowbreak \\\dfrac{\displaystyle 1}{\displaystyle 4},\allowbreak &0\le x<2,\allowbreak \\0,\allowbreak &\text{其他}.\end{cases}\) 令 \(\displaystyle Y=X^2\)，\(\displaystyle F(x,\allowbreak y)\) 为二维随机变量 \(\displaystyle (X,\allowbreak Y)\) 的分布函数。求：（I）\(\displaystyle Y\) 的概率密度 \(\displaystyle f_Y(y)\)；（II）\(\displaystyle \operatorname{Cov}(X,\allowbreak Y)\)；（III）\(\displaystyle F(-\dfrac{\displaystyle 1}{\displaystyle 2},\allowbreak 4)\)。


（23）（本题满分13分）设总体 \(\displaystyle X\) 的概率密度为 \(\displaystyle f(x;\theta)=\begin{cases}\theta,\allowbreak &0<x<1,\allowbreak \\1-\theta,\allowbreak &1\le x<2,\allowbreak \\0,\allowbreak &\text{其他},\allowbreak \end{cases}\) 其中 \(\displaystyle \theta\) 是未知参数 \(\displaystyle (0<\theta<1)\)，\(\displaystyle X_1,\allowbreak X_2,\allowbreak \ldots,\allowbreak X_n\) 为来自总体 \(\displaystyle X\) 的简单随机样本，记 \(\displaystyle N\) 为样本值 \(\displaystyle x_1,\allowbreak x_2,\allowbreak \ldots,\allowbreak x_n\) 中小于1的个数。求：（I）\(\displaystyle \theta\) 的矩估计；（II）\(\displaystyle \theta\) 的最大似然估计。

\end{document}
